\documentclass[a4paper,12pt]{article}
\usepackage[slovene]{babel}
\usepackage[utf8]{inputenc}
\usepackage[T1]{fontenc}
\usepackage{lmodern}
\usepackage{amsthm}
\usepackage{url}

{\theoremstyle{definition}
\newtheorem{definicija}{Definicija}}
{\theoremstyle{plain}
\newtheorem{izrek}{Izrek}}

% Definicija okoli za definicije in vaje
% Nastavimo koliko prostora izpustimo pod napisom nad sliko ali tabelo.
\setlength{\belowcaptionskip}{2mm}
% Nastavimo odmik prve vrstice v odstavku na 0mm.
\setlength{\parindent}{0mm}

% Kjer v besedilu naletite na niz `!!`, ga nadomestite z matematičnim izrazom
% iz datoteke `magic-resitev.pdf`.

%%%%%%%%%%%%%%%%%%%%%%%%%%%%%%%%%%%%%%%%%%%%%%%%%%%%%%%%%%%%%%%%%%%%%%%%
% Prvi sklop nalog
%%%%%%%%%%%%%%%%%%%%%%%%%%%%%%%%%%%%%%%%%%%%%%%%%%%%%%%%%%%%%%%%%%%%%%%%

\title{Magični kvadrati}
\date{}
% 1. naloga
% Glavo dokumenta popravite tako, da se ne bo izpisal datum.
% To naredite tako, da za ukazom `\title` dodate ukaz `\date{}`.

\begin{document}

\maketitle

Prirejeno po:
%
% 2. naloga
% Za izpis spletnih naslovov morate uporabiti ukaz `\url`.
% S tem ukazom lahko izpišete en spletni naslov tako, da
% naslov podate kot argument ukaza `\url` v zavitih oklepajih.
% V spodnjem okolju `itemize` uporabite ukaz `\url` na obeh spletnih naslovih
% (naslova sta v komentarju, da ne pride do napake pri prevajanju te datoteke).
\begin{itemize}
   \item \url{http://mathworld.wolfram.com/MagicSquare.html}
   \item \url{http://en.wikipedia.org/wiki/Magic_square}
\end{itemize}

% 3. naloga
% Na tem mestu uporabite ukaz, ki izpiše kazalo vsebine.
% Poiščite ga na spletu: po angleško se kazalu vsebine reče "table of contents".
\tableofcontents

%%%%%%%%%%%%%%%%%%%%%%%%%%%%%%%%%%%%%%%%%%%%%%%%%%%%%%%%%%%%%%%%%%%%%%%%
% Konec prvega sklopa nalog
%%%%%%%%%%%%%%%%%%%%%%%%%%%%%%%%%%%%%%%%%%%%%%%%%%%%%%%%%%%%%%%%%%%%%%%%

% Z ukazom `\newpage` lahko povzročite prehod na novo stran.
\newpage

\section{Uvod}

 \begin{definicija}
   \emph{Magični kvadrat} reda \(n\) je nabor \(n^2\) različnih števil,
   ki so razvrščena v kvadratno tabelo tako, da vedno dobimo enako vsoto,
   če seštejemo vsa števila poljubne vrstice, vsa števila poljubnega
   stolpca ali vsa števila v katerikoli od glavnih diagonal.
 \end{definicija}

Primer magičnega kvadrata reda 3 je prikazan v tabeli ??. % sklic: table:mag3

   
      
\begin{table}[h!]
   
\caption{Magični kvadrat reda}
\label{table:mag3}
% - napis nad tabelo: Magični kvadrat reda 3
% - oznaka: table:mag3
% Začetek magičnega kvadrata
\begin{center}
\begin{tabular}{|c|c|c|}
   \hline
   8 & 1 & 6 \\\hline
   3 & 5 & 7 \\\hline
   4 & 9 & 2 \\
   \hline
\end{tabular}
\end{center}

\end{table}
% Konec magičnega kvadrata


 \begin{definicija}
   Magični kvadrat reda \(n\) je \emph{normalen}, če v njem nastopajo števila
   \begin{equation}
      % oznaka: eq:numbers
      1, 2, 3, \ldots, n^2-1, n^2.
   \end{equation}
 \end{definicija}

Magični kvadrat v tabeli ?? je normalen.
To je tudi najmanjši netrivialen normalen magični kvadrat.
Poleg normalnih magičnih kvadratov so zanimivi tudi magični kvadrati praštevil. % sklic: table:mag3

%%%%%%%%%%%%%%%%%%%%%%%%%%%%%%%%%%%%%%%%%%%%%%%%%%%%%%%%%%%%%%%%%%%%%%%%

\section{Zgodovina}

\subsection{Kvadrat ">Lo Shu"<}

Kitajska literatura iz časa vsaj 2800 let pred našim štetjem govori o legendi
\emph{Lo Shu} -- ">zvitek reke Lo"<. V antični Kitajski je prišlo do
silne poplave. Ljudje so skušali rečnemu bogu narasle reke Lo ponuditi daritev,
da bi pomirili njegovo jezo. Iz vode se je prikazala želva z zanimivim vzorcem
na oklepu: v tabeli velikosti tri krat tri so bila predstavljena števila, tako
da je bila vsota števil v katerikoli vrstici, kateremkoli stolpcu in na obeh
glavnih diagonalah enaka: 15. To število je tudi enako številu dni v 24 ciklih
kitajskega sončnega leta. Ta vzorec so na določen način uporabljali upravljalci
reke.

%%%%%%%%%%%%%%%%%%%%%%%%%%%%%%%%%%%%%%%%%%%%%%%%%%%%%%%%%%%%%%%%%%%%%%%%

\subsection{Kulturna pomembnost}

Magični kvadrati so fascinirali človeštvo skozi vso zgodovino. Najdemo jih
v številnih kulturah, npr.\ v Egiptu in Indiji, vklesane v kamen ali
kovino, uporabljane kot talismane za dolgo življensko dobo in v
izogib boleznim.

\emph{Kubera-Kolam} je talna poslikava, ki se uporablja v Indiji, in je v
obliki magičnega kvadrata reda 3. Ta je v bistvu enak kot kvadrat
Lo Shu, vendar je vsako število povečano za 19.

Z magičnimi kvadrati so se ukvarjali tudi najbolj znani matematiki kot na
primer Euler, glej \cite{euler}

%%%%%%%%%%%%%%%%%%%%%%%%%%%%%%%%%%%%%%%%%%%%%%%%%%%%%%%%%%%%%%%%%%%%%%%%

\subsection{Zgodnji kvadrati reda 4}

Najzgodnejši znani magični kvadrat reda 4 je bil odkrit na napisu
v Khajurahu v Indiji in v Enciklopediji Bratovščine Čistosti iz enajstega
ali dvanajstega stoletja. Vrh vsega gre celo za ">panmagični kvadrat"<.
V Evropi sta morda najbolj znana naslednja magična kvadrata reda 4.

\paragraph{Magični kvadrat v litografiji Melancholia I}
Albrechta Dürerja naj bi bil najzgodnejši magični kvadrat
v evropski umetnosti. Zelo podoben je kvadratu Yang Huija, ki je nastal na Kitajskem
približno 250 let pred Dürerjevim časom.

Vsoto 34 je mogoče najti pri seštevanju števil v vsaki vrstici, vsakem stolpcu,
na vsaki diagonali, v vsakem od štirih kvadrantov, v sredinskih štirih poljih,
v štirih kotih, v štirih sosedih kotov v smeri urinega kazalca (\(3+8+14+9\)), v
štirih sosedih kotov v nasprotni smeri urinega kazalca (\(2+5+15+12\)), v dveh naborih
simetričnih parov (\(2+8+9+15\) in \(3+5+12+14\)), in še na nekaj drugih načinov.
Števili na sredini spodnje vrstici tvorita letnico litografije: 1514.

\paragraph{Pasijonska fasada na katedrali Sagrada família v Barceloni}
vsebuje magični kvadrat reda 4.

Vsota števil v vrsticah, stolpcih oziroma na diagonalah je 33 -- Jezusova starost
v času pasijona. Strukturno je kvadrat podoben Dürerjevemu, vendar so števila
v štirih poljih zmanjšana za 1. Posledica je, da sta števili 10 in 14 podvojeni
in zato kvadrat ni normalen.

%%%%%%%%%%%%%%%%%%%%%%%%%%%%%%%%%%%%%%%%%%%%%%%%%%%%%%%%%%%%%%%%%%%%%%%%

\section{Osnovne lastnosti}

\begin{definicija}
      Vsoto ene vrstice, enega stolpca ali ene od glavnih diagonal
      v magičnem kvadratu imenujemo \emph{magična konstanta}.
\end{definicija}

% Začetek okolja izrek
   
   \begin{izrek}
      Magična konstanta normalnega magičnega kvadrata reda !!
      je enaka
      \begin{equation}
         % oznaka: eq:mc
         % Okolje `equation` je že prikazni matematični način,
         % zato tu izraz vpišite brez `\[` in `\]`.
         !!
      \end{equation}
   \end{izrek}
% Konec okolja izrek

% Začetek dokaza
   
   \begin{proof}
      V normalnem magičnem kvadratu reda \(n\) je vsota vseh nastopajočih
      števil (glej ?? na strani ??) enaka
      !!. Ker imamo
      v kvadratu \(n\) vrstic z enako vsoto, je vsota števil v eni vrstici
      enaka številu \(M_2(n)\). % sklica: eq:numbers, eq:numbers
   \end{proof}
% Konec dokaza

Preprost račun pokaže, da je konstanti ?? analogna konstanta
\(M_2(n;A,D)\) za magični kvadrat, v katerem so nameščena števila
\(A\), \(A+D\), \(A+2D\), \ldots, \(A+(n^2-1)D\), enaka % sklic: eq:mc
!!

\begin{definicija}
      Če vsako od števil v normalnem magičnem kvadratu reda \(n\) odštejemo
      od števila !!, dobimo nov magični kvadrat, ki je prvotnemu
      \emph{komplementaren}.
\end{definicija}

\begin{definicija}
      Pravimo, da sta dva magična kvadrata \emph{različna}, če enega ni mogoče dobiti
      iz drugega s pomočjo zasukov oziroma zrcaljenj.
\end{definicija}

Števila različnih normalnih magičnih kvadratov se nahajajo v tabeli ??.
% sklic: table:stevila

% - napis nad tabelo: Število različnih normalnih magičnih kvadratov
% - oznaka: table:stevila
% Začetek tabele
      točna vrednost približek
      red 1 2 3 4 5 6
      % Decimalno vejico v matematičnem načinu zapišite kot {,},
      % sicer jo LaTeX izpiše kot ločilo, s presledkom za njo.
      število kvadratov 1 0 1 880 275305224 !!
% Konec tabele

Vse normalne magične kvadrate reda 4 je oštevilčil Frénicle de Bessy
leta 1693, glej \cite{bessy}, in jih je moč najti v knjigi \cite{berlekamp}
iz leta 1982. Število normalnih kvadratov reda 5 je izračunal
R. Schroeppel leta 1973 (glej Gardner \cite{gardner}).
Natančno število vseh različnih normalnih magičnih kvadratov reda 6 ni znano.
Avtorja navedenega približka sta Pinn in Wieczerkowski (glej \cite{pinn}), ki
sta za oceno uporabila simulacijo Monte Carlo in metode statistične mehanike.
 % ključi: bessy, berlekamp, gardner, pinn

%%%%%%%%%%%%%%%%%%%%%%%%%%%%%%%%%%%%%%%%%%%%%%%%%%%%%%%%%%%%%%%%%%%%%%%%

\newpage
\bibliographystyle{siam}
\bibliography{magic}
% Tu vstavite bibliografijo

%%%%%%%%%%%%%%%%%%%%%%%%%%%%%%%%%%%%%%%%%%%%%%%%%%%%%%%%%%%%%%%%%%%%%%%%

\end{document}
